\chapter{Conclusion and Future Work}
\label{ch:conclusion}

With the related systems in view, we return to the research questions,
describe future work, and close with what the work has shown and where it
stands in relation to high-performance computing.

\section{Answers to the Research Questions}
\label{sec:answers}

We asked whether treating the traffic of split inference as mixed, that is,
as small urgent messages and large patient ones, helps on commodity hardware.
The three research questions that made this precise are repeated here as
posed, each with its answer, the evidence for it, and its limit.

The first question concerned the separation itself. \emph{\rqone}
\waiting{results from E3 and E4: state the answer for each stack in one
sentence and fill the blanks} Part of the evidence needs no measurement,
since with one endpoint per class no cold message is ever ahead of a hot
message in the send queue of the hot class (\cref{sec:pool}). The rest is
measured, as the packets of both endpoints still meet on one device and one
cable. On the lab
pair, the 99th percentile of a 64-byte round trip under a stream of 1 MiB
writes is \blank\us{} with one endpoint per class and \blank\us{} on a shared
endpoint, against \blank\us{} with no stream, and on the home pair the three
values are \blank, \blank{} and \blank\us{} (\cref{tab:e3,tab:e4}). The
answer is limited to
one pair of machines on one cable, with one sender and without a switch or
packet loss, and it is an answer for the pool: on the serving path the sender
of the connector still serializes both classes in one queue
(\cref{sec:connector}).

The second question concerned the price of portability. \emph{\rqtwo}
\waiting{results from E1, E2 and E5: state the answer in one sentence and
fill the blanks} The comparison is a controlled one, as the two implementations stand behind
one interface and are measured by the same benchmark source
(\cref{sec:pool,sec:sfbench}). On the lab pair, a 64-byte write takes
\blank\us{} at the median on raw verbs and \blank\us{} on libfabric with
connected endpoints, and the cold class reaches \blank{} and \blank{} MB/s on
the two stacks for 1 MB writes (\cref{tab:e1,tab:e2}), while \cref{tab:e5}
gives the cost of each completion signal in instructions per operation. The
answer holds for one version of the library and for its verbs provider on
these two pairs of machines.

The third question concerned the request as a user sees it. \emph{\rqthree}
Arithmetic bounds the transfer at no less than 134 ms for a 2,048-token prompt on the laboratory link
(\cref{eq:wiretime}), and for that prompt the measured transfer takes
\blank{} ms, which is \blank{} percent of the time to first token
(\cref{tab:e8split}). \waiting{results from E8: state the answer to the
first half and fill the blanks} Its second half is answered for
each component separately and not for the three together, as the components
that decide on routing, compression and spilling are evaluated outside the
serving path (\cref{sec:overview}). The quantizer halves the number of bytes
at a measured cost in attention agreement (\cref{sec:res-quant}), host
staging is a small share of the wire time at this link speed under an
assumed copy rate (\cref{eq:staging}), and the storage tier is measured on
its own (\cref{sec:res-nvme}). This answer has the widest limit, as it
rests on one model, on random prompts, and on a serving path in which the
class of a transfer follows from its size.

\section{Future Work}

The work can be extended in four directions, given here from the nearest to
the furthest. The nearest is to complete the serving path. The decision of the rule tree
can be passed to the sender as the class of each blob and to the storage
tier as its spill decision, so that E8 can compare routing, compression and
spilling against the size-based path measured first (\cref{sec:overview}).
The sender of the connector serializes both classes in one queue, and a
queue per class would carry the separation of the pool up to the connector
(\cref{sec:connector}).

The second direction is to repair and tune the components that decide. The
switching rule of the controller should be replaced with the
routing-consistent one, or the controller turned into a service allocator,
and then measured. The classifier should be retrained on the signal log from
the cluster, labeled by measured latency, and the thresholds of the default
policy, which are initial values, should be set from measured traffic. The
sketch should be either used or removed: either the receiving side applies
the correction, or the sender stops sending 32 bytes per vector that nothing
reads, with the quality of each choice measured on generated text.

The third direction is the hardware. On server 4, installing the rebuilt
kernel and loading the peer memory module would open the GPU-direct path,
and the cost of staging could then be measured by removing it on the one
machine where that is possible. Beyond one pair of machines, the design has
to be tested with a switch, with several senders and with packet loss, and a
parallel file system reached over the same transport would add a tier below
local flash.

The furthest direction is the next transport. The pool can be built against
libfabric 2.x with per-endpoint memory registration, and, once an Ultra
Ethernet device is available, it can be tested whether separate endpoints
help over a transport that spreads packets across paths
(\cref{sec:uec-port}).

\section{Conclusion}

The limit on computing has moved from arithmetic, to memory, to the path
between machines, and high-performance computing is the industry that grew
up around the third limit. Its largest machines come from a handful of
builders, most of their interconnects come from two companies, and their
networks are engineered to keep one kind of traffic from delaying another
(\cref{sec:history}). Language-model inference reached the same limit within
a few years, and most work on it assumes the network of a data center.

This work asked what can be done without any of that. The laboratory pair
has no RDMA hardware, three of the four servers carry consumer GPUs that no
network card can reach directly, one board would not start with its
accelerators until its firmware and an ACPI table had been changed, and the
fastest network cards arrived configured for a different network. Despite these conditions, persistence prevailed. Each
of these obstacles was addressed in turn (\cref{ch:bringup}), with the
GPU-direct path the one that remains closed, and on machines that were never meant for the task we built
in software what larger systems obtain from dedicated hardware: a transport
with separate traffic classes on two libraries behind one interface and a
flow-controlled channel whose properties are proved, together with a
connector into a production serving engine and a method in which every
number carries the commands that produced it.

Performance followed. In the bytes that are sent, the quantizer halves a
real cache relative to 16-bit floating point while the eight most-attended
positions agree 91 percent of the time, and batching raises the write
throughput of the storage tier by 25 percent. In the order in which they
are sent, a 64-byte message under a stream of 1 MiB writes has a 99th
percentile of \blank\us{} with one endpoint per class against \blank\us{} on
a shared endpoint on the laboratory pair, and the cache of a 2,048-token
prompt, which cannot cross the laboratory link in less than 134 ms, reaches
the decode server in \blank{} ms. \waiting{results from E3 and E8:
fill the blanks, and keep the opening sentence of this paragraph only if
they support it} The analysis also removed
two components from the list of those that can be trusted, as we proved that
a classifier built on size thresholds cannot recognize decode traffic and
showed that the switching rule of the controller contradicts its own
derivation.

High-performance computing is the reference for this work in both
directions. The design borrows what that field already practices, namely
one-sided access to remote memory, which Cray built into the T3D, and the
separation of traffic classes, which Slingshot provides inside the network.
In return, the middleware is written against libfabric, the interface
through which Slingshot is programmed and which Ultra Ethernet has adopted,
so a design that is proved on two ordinary servers and one cable speaks the
interface of the largest machines. For a laboratory with the same hardware
and the same question, the library, the benchmark, the proofs and the record
of what it took to make the machines work are meant to serve as a reference
in turn. The idea itself is simple: a message of eight bytes should not
wait behind a block of four megabytes. \waiting{results from E3: close with
what the software network card allowed}
